\chapter{Reproduction}\label{ch:reproduction}

\NewV\ \emph{The weakest chapter in the book, placed last and labelled as
such. It proposes a canonical morphism and claims nothing beyond that.}

\begin{workplan}{\NewV}
\wpgoal{Give ``the kernel reproduces itself'' a precise candidate definition,
and be explicit that a canonical morphism is not yet a viable offspring.}

\wpsource{\texttt{TATIANA\_V1\_IDENTITY\_MAP.md}~\S5 (the proposal and its own
warning), \S10 open decision~4;
\texttt{MATH\_TOOLBOX.md}~II.2 (the adjunction, which is standard and already
proved).}

\wpblocked{Chapter~\ref{ch:identity} (there is no sub-invariant to carry until
$\mathrm{K}(\mathbb{K})$ is established) and Chapter~\ref{ch:twophase} (there
is nothing for an offspring to \emph{do} until the wake/sleep cycle is
specified). Do not write this chapter before those two are finished.}

\wpsteps{
  \item \textbf{Say what reproduction is not.} Not a copy: a bitwise
    duplicate carries no information about what the kernel learned to be, and
    duplicating the whole store is not growth but storage. State this before
    proposing anything.
  \item \textbf{Propose the morphism.} The adjoint pair is already in the
    mathematics: $\iota^*$ instantiates a downward-closed subcomplex,
    $\iota_!$ writes back only what was touched, and because $\iota$ is a full
    subposet inclusion the unit is an isomorphism, $\iota^*\iota_! \cong \Id$
    --- so the round trip provably does not corrupt the fragment. An
    offspring is a downward-closed subcomplex together with its pullback
    sheaf.
  \item \textbf{State the invariant it carries.} The offspring's
    $\mathrm{K}$ is the restriction of the parent's dimension vector to the
    subcomplex --- a \emph{sub-invariant}. Establish that this is well-defined
    and that it is conserved by the offspring's own consolidation, which
    follows from Chapter~\ref{ch:identity} applied to the subcomplex.
  \item \textbf{Do not claim viability.} Tag the whole chapter \Proposed. The
    honest statement is: the morphism is canonical, and nothing follows from
    that about whether the offspring can run, learn, or survive contact with
    input. Say so in those words.
  \item \textbf{Write the falsifier.} Before any implementation: what
    observation would show an offspring is \emph{not} viable? Candidates to
    develop --- the offspring's wake phase never stalls (so it can never
    grow); or it stalls immediately everywhere (so its address carries no
    information); or its invariant drifts where the parent's did not. Specify
    at least one as a runnable experiment.
  \item \textbf{Connect to the selection question, and leave it open.} If
    reproduction is cheap and viability is rare, something must choose which
    fragments reproduce. Note that V0 explicitly \Refuted\ random
    variation-and-selection as the growth mechanism --- variation needs
    randomness \emph{because it has no address}, and this architecture
    computes an address --- so the same argument constrains what a selection
    rule here may look like. Do not propose one yet.
}

\wpdone{The morphism is defined, the sub-invariant is proved to be
well-defined, the viability claim is explicitly \emph{not} made, and at least
one falsifier is specified precisely enough to run.}
\end{workplan}
